
le composé de deux morphismes parfaits, et par suite un morphisme parfait.

\textbf{Proposition 1.10.} \uline{Soient} \(X\), \(Y\) \uline{deux schémas lisses sur} \(S\). \uline{Alors toute} \(S\)-\uline{immersion de} \(X\) \uline{dans} \(Y\) \uline{est régulière}. \uline{En particulier}, \uline{toute section d'un morphisme lisse est une immersion régulière}.

La proposition résulte immédiatement de EGA IV 17.12.1, puisqu'une immersion régulière au sens de EGA IV est régulière au sens de 1.4.

\section*{2. Calculs sur les immersions régulières}

\subsection*{2.1.}
Soient \(A\) un anneau, \(B\) et \(C\) deux \(A\)-algèbres. On se propose de munir \(\bigoplus_{i\ge0}\Tor_i^A(B,C)\) d'une structure de \(A\)-algèbre graduée, donc de définir des homomorphismes :
\begin{equation}\tag{2.1.1}
 \Tor_i^A(B,C)\otimes_A\Tor_j^A(B,C)\longrightarrow \Tor_{i+j}^A(B,C).
\end{equation}

Soit \(P_\bullet\) une résolution de \(B\) par des \(A\)-modules projectifs; les \(\Tor_i^A(B,C)\) sont les groupes d'homologie du complexe \(P_\bullet\otimes_A C\). On sait définir des homomorphismes canoniques ([1], IV 6):
\begin{equation}\tag{2.1.2}
 H_i(P_\bullet\otimes_A C)\otimes_A H_j(P_\bullet\otimes_A C)
 \longrightarrow H_{i+j}\bigl((P_\bullet\otimes_A C)\otimes_A(P_\bullet\otimes_A C)\bigr).
\end{equation}

Par ailleurs, \(P_\bullet\otimes_A P_\bullet\) est un complexe de \(A\)-modules projectifs, à degrés \(\geq0\), et \(H_0(P_\bullet\otimes_A P_\bullet)=B\otimes_A B\); de plus, \(P_\bullet\) est une résolution de \(B\). Par suite, le morphisme d'augmentation \(B\otimes_A B\to B\) se prolonge en un morphisme de complexes : \(P_\bullet\otimes_A P_\bullet\to P_\bullet\) et définit donc, compte tenu de l'homomorphisme \(C\otimes_A C\to C\) déduit de la structure d'algèbre de \(C\), des homomorphismes :
\begin{equation}\tag{2.1.3}
 H_{i+j}\bigl((P_\bullet\otimes_A C)\otimes_A(P_\bullet\otimes_A C)\bigr)
 \longrightarrow H_{i+j}(P_\bullet\otimes_A C).
\end{equation}
